\documentstyle[%


righttag%


,amssymb%


,11pt%


,russian%


,izvru%


]{amsart}





\newcommand{\tl}[3]{\noindent{\bf#1}\ #2 \dotfill #3 \par} 


\setcounter{page}{93}


\pagestyle{plain}


\begin{document}


\par


\vskip 2cm


\par


\bigskip


\bigskip


\bigskip


\par


{\centerline{\Large СОДЕРЖАНИЕ ЖУРНАЛА ``ИЗВЕСТИЯ ВЫСШИХ УЧЕБНЫХ ЗАВЕДЕНИЙ''}}





{\centerline{\Large серия ``Математика'' за 2000 г.}}


\bigskip


\bigskip








\tl{Абрашин В.Н., Егоров А.А., Жадаева Н.Г.}


{Экономичные итерационные схемы\newline реализации метода 


конечных элементов для стационарных краевых задач 


математической\newline физики}{\No\,11}


\tl{Азбелев Н.В., Симонов П.М.}{Устойчивость уравнений с 


запаздывающим\newline аргументом. II}{\No\,4}


\tl{Акишев Г.А., Махашев С.Т.} 


{Об абсолютной сходимости рядов Фурье 


по обобщенной\newline системе Хаара}{\No\,3}


\tl{Аксентьева Е.П.}{Степенная задача Римана с постоянными


показателями для двоякопери-\newline одических функций с нулями


на контуре}{\No\,10}


\tl{Аксенюшкина Е.В.}


{Итерационный метод решения линейных 


задач оптимального\newline управления с терминальными ограничениями}{\No\,12}


\tl{Александров А.Ю.}{К вопросу об устойчивости


по неавтономному первому приближению}{\No\,10}


\tl{Алексеев Д.В.} 


{Приближение функций нескольких переменных


с весом\newline Чебышева--Эрмита}{\No\,6}


\tl{Алексенко Н.В.}


{Устойчивость решений нелинейных почти периодических систем\newline


функционально-дифференциальных


уравнений запаздывающего типа}{\No\,2}


\tl{Аллаков И.А.}


{О представлении четных чисел суммой двух нечетных простых


чисел из\newline арифметической прогрессии}{\No\,8}


\tl{Амосов Г.Г.}


{Об аппроксимации полугрупп изометрий


в гильбертовом пространстве}{\No\,2}


\tl{Андронов В.Г., Белоусов Е.Г.} 


{Об алгебраических свойствах множеств,


не содержащих\newline прямых}{\No\,5}


\tl{Антонова О.В., Бронникова С.А.,


Давыдова В.В.,  Кабрера М.\,Ордоньез.}


{О слабом\newline законе больших чисел в банаховых


пространствах мартингального типа $p$ 


при общем условии\newline типа Чезаро}{\No\,3}


\tl{Арутюнянц Г.В.}


{Об ограниченности некоторых


максимальных и мультипликаторных\newline операторов}{\No\,3}


\tl{Ахтямов А.М.} 


{О коэффициентах разложений по собственным функциям  


краевой\newline задачи со спектральным параметром  


в граничных условиях}{\No\,2}


\tl{Белоусов Е.Г., Андронов В.Г.} 


{Об алгебраических свойствах множеств,


не содержащих\newline прямых}{\No\,5}


\tl{Бикчантаев И.А.}


{Оценка числа решений однородной задачи Римана


на параболической\newline римановой поверхности}{\No\,1}


\tl{Бикчантаев И.А.}


{Задача Римана на ультрагиперэллиптической поверхности}{\No\,2}


\tl{Бронникова С.А., Антонова О.В., 


Давыдова В.В., Кабрера М.\,Ордоньез.}


{О слабом\newline законе больших чисел в банаховых


пространствах мартингального типа $p$ 


при общем условии\newline типа Чезаро}{\No\,3}


\tl{Бугир М.К.} 


{Условия неколеблемости решений систем дифференциальных уравнений 


\newline второго порядка}{\No\,5}


\tl{Васильев О.В., Васильева О.О.}


{К поиску равновесных управлений в дифференциальной\newline игре


$m$ лиц}{\No\,12}


\tl{Васильева Е.В.}


{Об одной задаче классификации множества управлений для\newline


параболического уравнения}{\No\,1}


\tl{Васильева О.О., Васильев О.В.}


{К поиску равновесных управлений в дифференциальной\newline игре


$m$ лиц}{\No\,12}


\tl{Винник А.В., Лейко С.Г.} 


{Изопериметрические экстремали поворота на двумерных \newline


связных группах Ли с инвариантными римановыми метриками}{\No\,7}


\tl{Власов В.В.}{Об оценках решений дифференциально-разностных уравнений  


нейтрального\newline типа}{\No\,4}


\tl{Вознюк И.П.}


{Аппроксимационный алгоритм для задачи размещения на максимум\newline


с ограниченными объемами производств и поставок}{\No\,12}


\tl{Волкодавов В.Ф., Мансурова Е.Р.} 


{Краевая задача для частного вида уравнения\newline Эйлера--Дарбу с 


интегральными условиями и специальными условиями сопряжения на \newline


характеристике}{\No\,8}


\tl{Гимади Э.Х., Кайран Н.М., Сердюков А.И.}


{О разрешимости многоиндексной\newline аксиальной задачи о назначениях


на одноциклических подстановках}{\No\,12}


\tl{Глазырина Л.Л., Павлова М.Ф.}


{Теорема о единственности решения одной задачи\newline теории


совместного движения русловых и подземных вод}{\No\,11}


\tl{Глекова В.В., Дубровский В.В.} 


{Теорема о единственности решения обратных задач\newline


спектрального анализа}{\No\,3}


\tl{Горбацевич В.В.}


{О тривиальности натурального расслоения некоторых компактных\newline


однородных пространств}{\No\,1}


\tl{Горлов С.К., Новиков И.Я., Родин В.А.} 


{Коррекция полиномов Хаара, применяемых \newline


для сжатия графической информации}{\No\,7}


\tl{Григорян М.Г.} 


{О системах универсальности в $L^p$, $1\le p<2$}{\No\,5}


\tl{Гулин А.В., Шередина А.В.}


{Границы устойчивости разностных схем}{\No\,11}


\tl{Гусенкова А.А., Плещинский Н.Б.}{Комплексные потенциалы


с логарифмическими\newline особенностями в ядрах


для упругих тел с дефектом вдоль гладкой дуги}{\No\,10}


\tl{Давыдова В.В., Антонова О.В., Бронникова С.А.,


Кабрера М.\,Ордоньез.}


{О слабом\newline законе больших чисел в банаховых


пространствах мартингального типа $p$ 


при общем условии\newline типа Чезаро}{\No\,3}


\tl{Данеев А.В., Русанов В.А.} 


{Об одном классе сильных дифференциальных моделей  


над\newline счетным множеством динамических процессов  


конечного характера}{\No\,2}


\tl{Дегтяренко Н.А.} 


{Решение в замкнутой форме 


интегрального уравнения типа свертки\newline


в гиперэллиптическом случае}{\No\,1}


\tl{Денисов А.М., Лоренци А.}


{О нелинейном интегральном уравнении первого рода}{\No\,11}


\tl{Драгилев М.М.} 


{О свойствах Абеля и Фабри базисных разложений аналитических \newline


функций}{\No\,8}


\tl{Дубровский В.В., Глекова В.В.} 


{Теорема о единственности решения обратных задач\newline


спектрального анализа}{\No\,3}


\tl{Егоров А.А., Абрашин В.Н., Жадаева Н.Г.}


{Экономичные итерационные схемы\newline реализации метода 


конечных элементов для стационарных краевых задач 


математической\newline физики}{\No\,11}


\tl{Емеличев В.А., Степанишина Ю.В.}


{Квазиустойчивость векторной нелинейной\newline траекторной задачи


с паретовским принципом оптимальности}{\No\,12}


\tl{Ерзакова Н.А.} 


{Меры некомпактности в исследовании неравенств}{\No\,9}


\tl{Ермолаева Л.Б.}


{Об одной квадратурной формуле}{\No\,3}


\tl{Ершова Т.В.} 


{Асимптотические теоремы для модификаций многочленов, 


подобных\newline многочленам Бернштейна}{\No\,9}


\tl{Жадаева Н.Г., Абрашин В.Н., Егоров А.А.}


{Экономичные итерационные схемы\newline реализации метода 


конечных элементов для стационарных краевых задач 


математической\newline физики}{\No\,11}


\tl{Жукова А.А.}


{Проблема Харди--Литтлвуда}{\No\,2}


\tl{Заботин И.Я.}


{Об устойчивости алгоритмов безусловной


минимизации псевдовыпуклых\newline функций}{\No\,12}


\tl{Заботин Я.И., Фукин И.А.}


{Об одной модификации метода сдвига штрафов для 


задач\newline нелинейного программирования}{\No\,12}


\tl{Захаров В.Н.} 


{Принцип абсолютного экстремума для одного 


класса дифференциальных\newline уравнений третьего 


порядка в трехмерном пространстве и его 


применение}{\No\,3}


\tl{Золотухина Л.А.}


{Асимптотическое распределение числа


совпадений двумерной выборки\newline при


натуральном совмещении}{\No\,5}


\tl{Иванков П.Л.} 


{О вычислении постоянных, входящих в оценки  


линейных форм}{\No\,1}


\tl{Иванова М.В., Ушаков В.И.}


{Свойства функции Грина третьей смешанной задачи для\newline


параболического уравнения в нецилиндрической области}{\No\,10}


\tl{Игошин В.А.} 


{Пульверизационное моделирование


и точечные симметрии пульверизации}{\No\,5}


\tl{Игошин В.А.} 


{Пульверизационное моделирование и точечно-траекторные 


морфизмы\newline квазигеодезических потоков}{\No\,7}


\tl{Кабрера М.\,Ордоньез, Антонова О.В., Бронникова С.А.,


Давыдова В.В.}


{О слабом\newline законе больших чисел в банаховых


пространствах мартингального типа $p$ 


при общем условии\newline типа Чезаро}{\No\,3}


\tl{Кайран Н.М., Гимади Э.Х., Сердюков А.И.}


{О разрешимости многоиндексной\newline аксиальной задачи о назначениях


на одноциклических подстановках}{\No\,12}


\tl{Какичев В.А., Нгуен Суан Тхао} 


{Базисный аналог $H$-функции одной и двух\newline переменных}{\No\,8}


\tl{Кац Б.А.}{Интеграл Стильтьеса по фрактальному контуру


и некоторые его приложения}{\No\,10}


\tl{Кехайопулу Н., Цингелис М.} 


{Замечание о полурешеточных конгруэнциях


в\newline упорядоченных полугруппах}{\No\,2}


\tl{Кехайопулу Н., Цингелис М.} 


{Замечание о конгруэнциях в упорядоченных\newline полугруппах}{\No\,5}


\tl{Кириллов С.А., Кузнецов М.И., Чебочко Н.Г.} 


{О деформациях алгебры Ли типа $G_{2}$\newline характеристики три}{\No\,3}


\tl{Комаров А.В., Пенкин О.М., Покорный Ю.В.} 


{О спектре равномерной сетки из\newline струн}{\No\,4}


\tl{Коннов И.В.}


{Приближенные методы для прямо-двойственных 


вариационных неравенств\newline смешанного типа}{\No\,12}


\tl{Кононов С.Г.}


{Алгебры инвариантных аффиноров


однородных пространств вещественных\newline простых групп Ли


типа $A_l$ с регулярными подгруппами изотропии}{\No\,5}


\tl{Коняев Ю.А.}  


{Спектральный метод исследования устойчивости  


некоторых классов\newline неавтономных дифференциальных уравнений}{\No\,5}


\tl{Коняев Ю.А.} 


{Метод расщепления в 


теории регулярных и сингулярных возмущений}{\No\,6}


\tl{Копп О.А. Тронин С.Н.}


{Матричные линейные операды}{\No\,6}


\tl{Корнеев В.Г., Фиш Я.}


{Двухуровневые методы для


пространственных задач, основанные\newline на агрегации}{\No\,11}


\tl{Крукиер Л.А., Чикина Л.Г.} 


{Кососимметрические итерационные методы решения\newline


стационарных задач конвекции-диффузии}{\No\,11}


\tl{Кузнецов М.И., Кириллов С.А., Чебочко Н.Г.} 


{О деформациях алгебры Ли типа $G_{2}$\newline характеристики три}{\No\,3}


\tl{Левенштам В.Б.} 


{О методе усреднения для 


квазилинейных параболических уравнений\newline с 


быстро осциллирующими коэффициентами}{\No\,7}


\tl{Лейко С.Г., Винник А.В.} 


{Изопериметрические экстремали поворота на двумерных \newline


связных группах Ли с инвариантными римановыми метриками}{\No\,7}


\tl{Лернер Э.Ю.} 


{Об адельной формуле для гауссовых интегралов}{\No\,7}


\tl{Лоренци А., Денисов А.М.}


{О нелинейном интегральном уравнении первого рода}{\No\,11}


\tl{Лукашенко Т.П.}{О свойствах обобщенных систем


разложения, подобных ортогональным}{\No\,10}


\tl{Маликов А.И.} 


{Об устойчивости систем дифференциальных уравнений со 


случайными\newline структурными изменениями}{\No\,1}


\tl{Маликов А.И.}  


{Матричные системы дифференциальных уравнений


с условием\newline квазимонотонности}{\No\,8}


\tl{Мансурова Е.Р., Волкодавов В.Ф.} 


{Краевая задача для частного вида уравнения\newline Эйлера--Дарбу с 


интегральными условиями и специальными условиями сопряжения на \newline


характеристике}{\No\,8}


\tl{Марченко И.И.} 


{Об оценке Шиа для величины отклонения


мероморфной функции}{\No\,8}


\tl{Марюкова Н.Е.} 


{Об одной геометрической интерпретации 


уравнения $\sin$-Гордона}{\No\,7}


\tl{Махашев С.Т., Акишев Г.А.} 


{Об абсолютной сходимости рядов Фурье 


по обобщенной\newline системе Хаара}{\No\,3}


\tl{Миротин А.Р.}


{Преобразование Лапласа и некоммутативный вариант


теоремы\newline Пэли--Винера}{\No\,9}


\tl{Можей Н.П.}


{Однородные подмногообразия в четырехмерной 


аффинной и проективной\newline геометрии}{\No\,7}


\tl{Морозов С.Ф., Семенов А.В.} 


{Обобщенные кривые и необходимые условия  


разрывного\newline решения пространственной вариационной задачи}{\No\,9}


\tl{Мунембе Ж.\,С.\,П.}{К вопросу об асимптотическом поведении


решений системы\newline линейных функционально-дифференциальных уравнений с


запаздыванием и\newline периодическими параметрами}{\No\,4}


\tl{Мухин Ю.Н.} 


{О топологических группах со счетным множеством 


монотетичных подгрупп}{\No\,2}


\tl{Насибов Ф.Г.}{Об экстремальных задачах в классах целых  


функций, ограниченных на\newline вещественной оси}{\No\,10}


\tl{Нгуен Суан Тхао, Какичев В.А.} 


{Базисный аналог $H$-функции одной и двух\newline переменных}{\No\,8}


\tl{Нежметдинов И.Р.}


{Двойственные дополнения и оболочки


классов аналитических\newline функций}{\No\,1}


\tl{Новиков И.Я., Горлов С.К., Родин В.А.} 


{Коррекция полиномов Хаара, применяемых \newline


для сжатия графической информации}{\No\,7}


\tl{Павлова М.Ф., Глазырина Л.Л.}


{Теорема о единственности решения одной задачи\newline теории


совместного движения русловых и подземных вод}{\No\,11}


\tl{Панюков А.В.}


{Экстремальный метод решения


параметрической обратной задачи


для\newline системы линейных функциональных уравнений}{\No\,9}


\tl{Пенкин О.М., Комаров А.В., Покорный Ю.В.} 


{О спектре равномерной сетки из\newline струн}{\No\,4}


\tl{Плещинский Н.Б.}


{К абстрактной теории приближенных методов 


решения линейных\newline операторных уравнений}{\No\,3}


\tl{Плещинский Н.Б., Гусенкова А.А.}{Комплексные потенциалы


с логарифмическими\newline особенностями в ядрах


для упругих тел с дефектом вдоль гладкой дуги}{\No\,10}


\tl{Покорный Ю.В., Комаров А.В., Пенкин О.М.} 


{О спектре равномерной сетки из\newline струн}{\No\,4}


\tl{Попов А.Ю., Теляковский С.А.}


{К оценкам частных сумм рядов Фурье функций\newline 


ограниченной вариации}{\No\,1}


\tl{Пудалова Е.И., Срочко В.А.}


{Методы нелокального улучшения допустимых\newline управлений


в линейных задачах с запаздыванием}{\No\,12}


\tl{Путилина А.В.}


{Асимптотика функции объема множества меньших


значений\newline полинома в $\Bbb R^n$}{\No\,6}


\tl{Ратыни А.К.} 


{Об эллиптической краевой задаче с оператором суперпозиции  


в граничном\newline условии. I}{\No\,4}


\tl{Рахматуллина Л.Ф.}


{Верхняя оценка спектрального радиуса изотонного оператора}{\No\,1}


\tl{Рахматуллина Л.Ф.}


{Непрерывная зависимость от параметров решения


линейной краевой\newline задачи}{\No\,4}


\tl{Родин В.А., Горлов С.К., Новиков И.Я.} 


{Коррекция полиномов Хаара, применяемых \newline


для сжатия графической информации}{\No\,7}


\tl{Русанов В.А., Данеев А.В.} 


{Об одном классе сильных дифференциальных моделей  


над\newline счетным множеством динамических процессов  


конечного характера}{\No\,2}


\tl{Рязанцева И.П.}


{Метод итеративной регуляризации второго порядка для


выпуклых\newline задач условной минимизации}{\No\,12}


\tl{Салимов Р.Б.} 


{Новый подход к решению краевой задачи Гильберта  


для аналитической\newline функции в многосвязной области}{\No\,2}


\tl{Сатаров Ж.С.}  


{Определяющие соотношения обобщенных ортогональных групп\newline над 


коммутативными локальными кольцами без единицы}{\No\,6}


\tl{Семенов А.В., Морозов С.Ф.} 


{Обобщенные кривые и необходимые условия  


разрывного\newline решения пространственной вариационной задачи}{\No\,9}


\tl{Семигродских А.П.}


{О клонах полиномов над бесконечными полями}{\No\,7}


\tl{Сердюков А.И., Гимади Э.Х., Кайран Н.М.}


{О разрешимости многоиндексной\newline аксиальной задачи о назначениях


на одноциклических подстановках}{\No\,12}


\tl{Серебряков В.П.} 


{Об индексе дефекта матричных дифференциальных операторов 


второго\newline порядка с быстро осциллирующими коэффициентами}{\No\,3}


\tl{Серебряков В.П.} 


{К вопросу об индексе дефекта матричных дифференциальных 


\newline операторов второго порядка}{\No\,5}


\tl{Сильченко Ю.Т.} 


{Об оценке резольвенты дифференциального оператора второго\newline


порядка с нерегулярными граничными условиями}{\No\,2}


\tl{Симонов П.М., Азбелев Н.В.}{Устойчивость уравнений с 


запаздывающим\newline аргументом. II}{\No\,4}


\tl{Сорокина М.М.} 


{О композиционных и локальных критических формациях}{\No\,7}


\tl{Сосов Е.Н.}


{Об $\varepsilon$-ядре ограниченного


множества в специальном метрическом пространстве}{\No\,9}


\tl{Срочко В.А., Пудалова Е.И.}


{Методы нелокального улучшения допустимых\newline управлений


в линейных задачах с запаздыванием}{\No\,12}


\tl{Степанишина Ю.В., Емеличев В.А.}


{Квазиустойчивость векторной нелинейной\newline траекторной задачи


с паретовским принципом оптимальности}{\No\,12}


\tl{Сумин М.И.} 


{Субоптимальное управление полулинейными  


эллиптическими уравнениями\newline с фазовыми  


ограничениями, I: принцип максимума для 


минимизирующих\newline последовательностей, 


нормальность}{\No\,6}


\tl{Сумин М.И.} 


{Субоптимальное управление полулинейными эллиптическими


уравнениями с\newline фазовыми ограничениями, II: чувствительность,


типичность регулярного принципа\newline максимума}{\No\,8}


\tl{Тер\"ехин М.Т.} 


{Об устойчивости управления по параметру}{\No\,9}


\tl{Тимербаев М.Р.}


{Конечноэлементная аппроксимация 


в весовых пространствах Соболева}{\No\,11}


\tl{Титов С.С.}


{Решение нелинейных уравнений 


в аналитических полиалгебрах. I}{\No\,1}


\tl{Титов С.С.}


{Решение нелинейных уравнений 


в аналитических полиалгебрах. II\ }{\No\,6}


\tl{Тонконог С.Л.} 


{Групповой анализ уравнений динамики неньютоновской жидкости}{\No\,9}


\tl{Тронин С.Н., Копп О.А.}


{Матричные линейные операды}{\No\,6}


\tl{Трынин А.Ю.} 


{Об отсутствии устойчивости интерполирования


по собственным функциям\newline задачи Штурма-Лиувилля}{\No\,9}


\tl{Ушаков В.И., Иванова М.В.}


{Свойства функции Грина третьей смешанной задачи для\newline


параболического уравнения в нецилиндрической области}{\No\,10}


\tl{\fbox{Фадеев Н.П.}} 


{Ортогональные многочлены обобщенно-четного веса}{\No\,8}


\tl{Федоров В.Е.} 


{Вырожденные сильно непрерывные


группы операторов}{\No\,3}


\tl{Фиш Я., Корнеев В.Г.}


{Двухуровневые методы для


пространственных задач, основанные\newline на агрегации}{\No\,11}


\tl{Фукин И.А., Заботин Я.И.}


{Об одной модификации метода сдвига штрафов для 


задач\newline нелинейного программирования}{\No\,12}


\tl{Цалюк З.Б.} 


{Асимптотическая структура резольвенты неустойчивого уравнения \newline


Вольтерра с разностным ядром}{\No\,4}


\tl{Цингелис М., Кехайопулу Н.} 


{Замечание о полурешеточных конгруэнциях


в\newline упорядоченных полугруппах}{\No\,2}


\tl{Цингелис М., Кехайопулу Н.} 


{Замечание о конгруэнциях в упорядоченных\newline полугруппах}{\No\,5}


\tl{Чебочко Н.Г., Кириллов С.А., Кузнецов М.И.} 


{О деформациях алгебры Ли типа $G_{2}$\newline характеристики три}{\No\,3}


\tl{Ченцов А.Г.}


{Универсальная версия обобщенных интегральных


ограничений в классе\newline конечно-аддитивных мер. II}{\No\,2}


\tl{Ченцов А.Г.} 


{К вопросу об итерационной реализации неупреждающих многозначных\newline


отображений}{\No\,3}


\tl{Чикина Л.Г., Крукиер Л.А.} 


{Кососимметрические итерационные методы решения\newline


стационарных задач конвекции-диффузии}{\No\,11}


\tl{Шагидуллин Р.Р.}


{Исследование положительных по кривизне


решений системы\newline уравнений равновесия замкнутой цилиндрической


оболочки}{\No\,11}


\tl{Шамоян Р.Ф.} 


{Непрерывные функционалы и мультипликаторы степенных рядов одного \newline


класса аналитических в поликруге функций}{\No\,7}


\tl{Шередина А.В., Гулин А.В.}


{Границы устойчивости разностных схем}{\No\,11}


\tl{Шипилева А.В.}


{Предельные распределения для ветвящихся процессов


с иммиграцией}{\No\,1}


\tl{Широкова Е.А.}


{Свед\'ение решения внутренней обратной краевой задачи


к\newline интегральному уравнению в случае угловых точек на искомом и


на известном контурах}{\No\,9}


\tl{Шпирко С.В.}


{Решение игр многих лиц с помощью градиентного метода прогнозного \newline


типа}{\No\,6}








\bigskip


{\centerline {КРАТКИЕ СООБЩЕНИЯ}}


\bigskip











\tl{Алвеш М.Ж.}


{Об одной нелинейной краевой задаче


с несуммируемой особенностью}{\No\,4}


\tl{Аминова А.В., Зуев С.В., Калинин Д.А.}


{Алгебраические условия согласования\newline двух метрических форм с общей


почти комплексной (кватернионной) структурой на\newline многообразии}{\No\,7}


\tl{Ахмадишина Ф.К.}


{Об асимптотической оптимальности 


проекционных методов}{\No\,3}


\tl{Аюпова Е.Ф. $($Мифтахова$)$} 


{О решении одного двумерного слабо сингулярного \newline


интегрального уравнения первого рода}{\No\,8}


\tl{Винник А.В., Лейко С.Г.} 


{Поворотно-конформные преобразования в  


плоскости\newline Лобачевского}{\No\,9}


\tl{Виноградова Г.Ю., Георгиев К.А., Деундяк В.М.}


{О структуре алгебр сингулярных\newline интегральных операторов со сдвигом


в весовых пространствах и вычисление индекса\newline семейств}{\No\,5}


\tl{Гайсин Т.И.}


{О комплексе базовых форм слоения}{\No\,7}


\tl{Ганиев И.Г., Чилин В.И.} 


{Индивидуальная эргодическая теорема для сжатий  \newline


в решетке $L_p(\wh\nabla,\wh\mu)$ Банаха--Канторовича}{\No\,7}


\tl{Георгиев К.А., Виноградова Г.Ю., Деундяк В.М.}


{О структуре алгебр сингулярных\newline интегральных операторов со сдвигом


в весовых пространствах и вычисление индекса\newline семейств}{\No\,5}


\tl{Гурова И.Н.}


{О методах полудискретизации для квазилинейных уравнений\newline


с некомпактной полугруппой}{\No\,4}


\tl{Гусаренко С.А.}


{О минимизации квадратичных функционалов


с ограничениями в виде\newline равенств}{\No\,4}


\tl{Денисюк И.Т.} 


{Одна задача сопряжения аналитических функций в аффинно \newline


преобразованных областях с кусочно-гладкими границами}{\No\,6}


\tl{Деундяк В.М., Виноградова Г.Ю., Георгиев К.А.}


{О структуре алгебр сингулярных\newline интегральных операторов со сдвигом


в весовых пространствах и вычисление индекса\newline семейств}{\No\,5}


\tl{Зубова С.П.}


{Исследование решения задачи Коши для одного


сингулярно возмущенного\newline дифференциального уравнения}{\No\,8}


\tl{Зуев С.В., Аминова А.В., Калинин Д.А.}


{Алгебраические условия согласования\newline двух метрических форм с общей


почти комплексной (кватернионной) структурой на\newline многообразии}{\No\,7}


\tl{Кадиев Р.И.}


{Достаточные условия устойчивости по части переменных


линейных\newline стохастических систем с последействием}{\No\,6}


\tl{Какичев В.А., Нгуен Суан Тхао.}


{Обобщенные свертки $H$-преобразований}{\No\,10}


\tl{Калинин Д.А., Аминова А.В., Зуев С.В.}


{Алгебраические условия согласования\newline двух метрических форм с общей


почти комплексной (кватернионной) структурой на\newline многообразии}{\No\,7}


\tl{Ковалев Л.В.}


{О внутренних радиусах


симметричных неналегающих областей}{\No\,6}


\tl{Кравцов В.М., Кравцов М.К., Лукшин Е.В.}


{О числе $r$-нецелочисленных


вершин\newline многогранника трехиндексной


аксиальной задачи выбора}{\No\,12}


\tl{Кравцов М.К., Кравцов В.М., Лукшин Е.В.}


{О числе $r$-нецелочисленных


вершин\newline многогранника трехиндексной


аксиальной задачи выбора}{\No\,12}


\tl{Крепкогорский В.Л.}


{Интерполяция в классе пространств Бесова и Лизоркина--Трибеля.\newline


Предельный случай $p_0=\infty$}{\No\,5}


\tl{Крючков А.Ф., Сиренек В.А.} 


{Вероятностное представление решения 


гиперболического\newline уравнения массопереноса 


с конвективным членом}{\No\,7}


\tl{Лейко С.Г., Винник А.В.} 


{Поворотно-конформные преобразования в  


плоскости\newline Лобачевского}{\No\,9}


\tl{Лукшин Е.В., Кравцов М.К., Кравцов В.М.}


{О числе $r$-нецелочисленных


вершин\newline многогранника трехиндексной


аксиальной задачи выбора}{\No\,12}


\tl{Молчанов А.В.} 


{Полугруппы эндоморфизмов слабых 


$p$-гиперграфов}{\No\,3}


\tl{Нгуен Суан Тхао, Какичев В.А.}


{Обобщенные свертки $H$-преобразований}{\No\,10}


\tl{Ожегова А.В.}


{О прямых методах решения одного класса


интегральных уравнений\newline первого рода}{\No\,5}


\tl{Ойнас И.Л., Цалюк З.Б.} 


{Асимптотика решений одной системы


разностных уравнений}{\No\,4}


\tl{Плехова Э.В.}


{Устойчивая разрешимость монотонных и аккретивных операторов}{\No\,4}


\tl{Сиренек В.А., Крючков А.Ф.} 


{Вероятностное представление решения 


гиперболического\newline уравнения массопереноса 


с конвективным членом}{\No\,7}


\tl{Сумин В.И., Чернов А.В.}


{О некоторых признаках квазинильпотентности\newline


функциональных операторов}{\No\,2}


\tl{Тимофеев Г.Н.} 


{Интегрируемая $\pi$-структура с симметричной  


второй ковариантной\newline производной аффинора}{\No\,5}


\tl{Федоров Д.Л.}


{О представлении резольвенты


интегральных уравнений Перрона--Стилтьеса\newline


при помощи рядов}{\No\,4}


\tl{Цалюк З.Б., Ойнас И.Л.} 


{Асимптотика решений одной системы


разностных уравнений}{\No\,4}


\tl{Чернов А.В., Сумин В.И.}


{О некоторых признаках квазинильпотентности\newline


функциональных операторов}{\No\,2}


\tl{Чилин В.И., Ганиев И.Г.} 


{Индивидуальная эргодическая теорема для сжатий  \newline


в решетке $L_p(\wh\nabla,\wh\mu)$ Банаха--Канторовича}{\No\,7}


\medskip











\noindent Рефераты статей, депонированных в ВИНИТИ 


через журнал ``Известия вузов.\newline Математика''\dotfill \No\,1








\end{document}


